\begin{pastexam}[label={px:ncku-109-comm-ds-04}]{成大電通所 109 學年度 第 4 題}
（是非；錯的要說明）

Consider the quick sort algorithm. True or false? Explain if false. (20 points)\\
(a) The average-case time complexity of this algorithm is the same as that of the merge sort
algorithm.\\
(b) It is stable.\\
(c) The worst-case time complexity of this algorithm is the same as that of the heap sort
algorithm.\\
(d) It is not suitable for external sorting.\\
(e) The best-case time complexity occurs when the input data are already sorted.
\end{pastexam}

\begin{pastsoln}
答案：(a) 對；(b) 錯；(c) 錯；(d) 對；(e) 錯。 以下的快速排序指一般的版本：取第一個（或最後一個）當
基準（資結書第~14.5~節）。

\textbf{(a) 對}。 快速排序平均 $\Theta(n \lg n)$（資結書定理~14.3），合併排序最好、平均、最壞都是
$\Theta(n \lg n)$。 兩者同階。

\textbf{(b) 錯}。 分割時會把元素跨過一大段交換，相同鍵值的元素相對順序可能被打亂，不穩定。 例子：用
Lomuto 分割（基準取最後一個）排 $\langle 2_a, 2_b, 1\rangle$（下標只是標記哪一個 $2$）：基準是 $1$，
$2_a$、$2_b$ 都比 $1$ 大，留在原處，最後基準和第一個大於它的 $2_a$ 交換，得到
$\langle 1, 2_b, 2_a\rangle$，右段 $\langle 2_b, 2_a\rangle$ 已經排好不會再換回來，結果 $2_b$ 跑到
$2_a$ 前面。

\textbf{(c) 錯}。 快速排序最壞 $\Theta(n^2)$（每次基準都是最大或最小，切成 $0$ 和 $m - 1$），堆積排序
最壞 $\Theta(n \lg n)$（資結書第~14.6~節）。 兩者不同。

\textbf{(d) 對}。 外部排序的資料放在磁碟上，只能大塊、循序地讀寫。 快速排序的分割要從兩端同時往中間
掃、到處交換，等於在磁碟上隨機存取，而且要整個陣列能隨機存取才能做；外部排序用的是只需要循序讀寫
的合併排序（先產生排好的段，再做 $k$ 路合併，資結書第~15.3~節）。

\textbf{(e) 錯}。 取第一個當基準時，已經排好的資料正好是\textbf{最壞}情況：基準永遠是最小值，每次切成
$0$ 和 $m - 1$，$\Theta(n^2)$。 最好的情況是每次基準都剛好是中位數，切成兩半，
$T(n) = 2T(n/2) + \Theta(n) = \Theta(n \lg n)$，和輸入是否排好無關。
\end{pastsoln}
